\subsection{联合迭代的计算过程}
\subsubsection{从完整方案到一次可接受的改进}
第~\ref{sec:q1} 节给出了方案空间、全局下界和接受准则。本节把候选从生成到保存的计算顺序具体化。对一个固定的计算图、问题编号和核数，搜索状态包含已验证的当前最好方案 $x$、对应完整执行结果、父方案代号、已经评价或安全剪枝的完整方案集合，以及各候选来源的近期记录。这里“当前最好”只有一个；不同来源产生的候选都与这个状态比较，来源自身不保存另一套最终解。初始构造可以来自本配置下已经验证的完整方案，或由依赖关键性、阶段粒度与负载映射生成。对初始候选先做结构检查，再取得完整执行结果；只有合法结果才进入当前最好状态。

一次迭代可按下列顺序表述，其中 $U=C_q(x)$，$L(y)$ 是固定候选 $y$ 的有效周期下界：
\begin{enumerate}
  \item 根据仍可用的来源、近期改善幅度、生成与评价成本、失败比例及剩余时间选择一个来源。初次覆盖和久未访问的来源获得机会；预计成本超过剩余时间时停止准入。
  \item 来源提出完整候选 $y=(\mathcal P',a',\boldsymbol\pi')$。检查非边界算子覆盖、每个子图恰有一次归核及商图与核内顺序合图的无环性；相同完整方案不重复评价。
  \item 若 $L(y)>U$，该候选不可能改善主目标，直接筛去。等号情形保留，因为周期持平时额外 COPY 字节仍可能下降。
  \item 对其余候选调用完整执行，重建任务和边界搬运、自动换出、Pipe 排布与共享带宽事件。只有结果可行，且 $(C_q(y),D_q(y))<_{\rm lex}(C_q(x),D_q(x))$，才令 $x\leftarrow y$ 并更新父方案代号；取得的解在评价和保存完成后才计为可用。
  \item 检查全局下界相对差距、评价次数、墙钟预算、来源耗尽和停滞条件。第三问结束主搜索后，对已选同一方案另作无缓存评价，配对结果不参与主搜索的接受判断。
\end{enumerate}
选择来源的反馈指数见式~\eqref{eq:sourcepriority}。实际准入先考虑少量来源预热与时间可行性，再比较近期收益和成本；它是调配有限评价机会的规则，不是候选周期的替代预测。每次接受都会改变父方案，故候选生成必须明确绑定接受前的状态。例如驻留候选使用先前的张量生命周期，缓存候选使用先前的读入与淘汰事件；若父方案改变，继续消耗旧候选流可能把旧事件错误关联到新划分或新归核。实现因此在新父方案代号出现后重建依赖父方案的结构、驻留和缓存流，并以已接受的完整方案及执行记录为起点。与父方案无关的全局构造流则保存其迭代位置，避免每次改进后重新生成相同初值。这个区分也解释了为何来源调用数不能直接看作互不干扰的独立实验。

\subsubsection{邻域更新、结构预检与计算量}
令 $N=|V|$ 为需归组的非边界算子数，$M$ 为这些算子间保留的依赖边数，$m$ 为当前子图数，$H(y)$ 为候选 $y$ 的完整执行实际处理的事件数。一次“迁移”把一个完整子图从原核心序列删除，再插入目标核心序列的某个合法位置；它更新 $a$ 与相关的 $\pi_k$，但不改变该子图内部算子集合。“交换”改变两个子图在一个或两个核心序列中的相对位置；若跨核交换，还同时更新两者归属。“拆分”把某子图分成非空、互不相交的两部分，在算子到子图映射中为新部分赋编号，并重新安置两部分的核心位置与先后顺序。联合迁移及张量关联移动按依赖闭合的子图集合操作，其输出仍须含全部 $N$ 个算子的完整映射和各核序列。拆分后的新增边界搬运、迁移后的跨核通信以及顺序变化造成的驻留区间，均由新方案的执行结果计算，不能从旧目标值直接加减。

给定一个已经构造好的完整候选，扫描 $N$ 个算子建立组成员与归属，扫描 $M$ 条依赖建立跨组前驱，再加入各核相邻子图的顺序边。后者总数不超过 $m$；在合图上做一次拓扑检查，按通常邻接表实现，其工作量为 $O(N+M+m)$，存储量同阶。这只是必要的结构预检复杂度，不是完整求解复杂度。环检查能够排除“某组必须先于另一组、又被核内顺序要求后于另一组”的矛盾，却不能决定瞬时缓冲是否可分配、自动换出是否引入额外请求、任务等待如何落在时间轴上，以及共享带宽在并发请求间如何分配。因此，通过预检的候选仍须完整执行验证。对事件评价记其实际成本为 $T_{\rm eval}(N,M,m,H(y))$；仅因事件轨迹含 $H(y)$ 项，不能未经实现级证明便称该成本对 $H$ 或 $N$ 线性。

候选空间本身也远大于一次预检。只枚举单个子图向其他核心迁移，已有至多 $m(n-1)$ 种归核选择；再考虑每个目标核心的插入位置、两个子图交换、组内依赖合法的拆分、多个关联组联动，候选数会继续增加。完整执行还有固定方案生成的 COPY、Spill 和 FIFO 事件，$H(y)$ 由候选实际轨迹决定，不能预先用原图边数替代。实际程序对单次来源流设生成上限，整配置最多评价 $2048$ 次，并以 $600$ 秒总预算及剩余时间准入约束搜索。若实际完成 $B$ 次评价，总代价应写成各次生成、结构检查、保存及 $\sum_{i=1}^{B}T_{\rm eval}(N,M,m_i,H_i)$ 的和，而非一个脱离候选数与轨迹长度的单一线性式。搜索中的下界剪枝、去重和早停减少实际调用，却不把有限预算算法变为全局枚举。

\subsubsection{两张五核图的方案展开}
真实五核方案能够说明“子图数”与“执行任务数”在两问中的差别。图80在第一问选出 $3$ 个子图，核心序列长度为 $(1,1,1,0,0)$：五个核心可用，但只有三个核心承担工作，三个子图分别构成 Task，最终总体任务执行时间为 $100325$ 周期，额外 COPY 为 $3047424$ B。第二问对同图独立寻优后得到 $161$ 个子图，五核序列长度为 $(34,32,34,30,31)$，周期为 $71369$，额外 COPY 为 $2273280$ B。第二问仍有许多子图，却只有五个核心级 Task；同核子图可延续存储状态，跨核交接按写出完成后等待的规则执行。由于两问同时改变了划分、归核和顺序，这组周期差只说明各自选中方案的表现，不能当作固定方案下删除任务等待的净效果。

图46的第一问五核方案含 $15$ 个子图，序列长度为 $(4,2,2,2,5)$。每个子图是独立 Task，所以五核序列中共有 $15-5=10$ 处同核相邻 Task 交接；承担五个子图的核心至少经历四次按题设定义的相邻任务等待。该方案周期为 $91730$，额外 COPY $3669568$ B，其中划分搬运 $3564096$ B、自动换出 $105472$ B。第二问独立选出的方案含 $31$ 个子图，序列长度为 $(5,7,7,5,7)$，周期 $69268$，额外 COPY $2469536$ B，未产生自动换出的额外字节。虽然第二问有更多子图，其同核子图不构成新的 Task 边界；第一问同核相邻 Task 的等待也不能简单乘以十再从 $91730$ 中相减，因为五核并行执行、跨核依赖和 DDR 服务会使这些等待与其他事件在时间上交错。两例分别展示了有效核心数、子图粒度和自动换出如何与任务规则共同进入完整周期。

\subsubsection{可行方案空间的嵌入性质}
现有下界说明任一可行结果距离全局最优最多仍可能改善多少。另有一项与核数有关、但不依赖搜索启发式的性质。固定问题编号和全部硬件参数，设允许核心空闲，增加一颗核心不减少原核心的 Pipe、L1/UB 容量，也不改变共享 DDR 或 L2 总带宽与事件优先规则。对任意 $n$ 核可行方案 $x$，构造 $n+1$ 核方案 $\iota(x)$：保持全部子图集合、算子映射、原有核心归属及原核序列不变，仅令新增核心的序列为空。由于没有任何新操作、COPY 或缓存请求从空核心发射，原方案的任务边界、缓冲实例和事件顺序均不变，故
\begin{equation}
 \iota(x)\in\mathcal F_q(n+1),\qquad
 C_q(\iota(x))=C_q(x),\qquad D_q(\iota(x))=D_q(x).
 \label{eq:core_embedding}
\end{equation}
于是同配置、同评价规则下的理论最优周期满足 $C_{q,n+1}^{*}\le C_{q,n}^{*}$。证明只需把任意 $n$ 核可行方案按上法嵌入 $n+1$ 核可行集；后者的最小值不大于其中这个方案的周期。该命题不要求新增核心必须被使用，也不意味着增加核心必然严格缩短周期。实际有限预算算法分别搜索不同核数，其已找到的两个方案并不构成包含关系；来源准入、候选顺序和搜索截止可能不同，所以观察到的近似结果不能单凭式~\eqref{eq:core_embedding} 声称逐图随核数严格单调。第一、二问的一核点还采用题设约定的原始整图参考，跨这一特殊基准解释时应使用实际统一评价定义，而不是把约定点误作上述嵌入构造的搜索输出。
